\section{Conclusiones y Trabajo Futuro}
\label{sec:conclusiones}

\subsection{Conclusiones}

La presente investigación ha implementado y validado un ecosistema integral de minería de datos, procesamiento de lenguaje natural y telemetría analítica multimodal sobre \nMensajes{} mensajes del canal oficial de Telegram de la Carrera de Informática en la Universidad Mayor de San Andrés (UMSA). Las conclusiones fundamentales del estudio son:

\begin{enumerate}[leftmargin=1.6em]
  \item \textbf{Prevalencia Factual en la Interacción Académica:} El \pctIndeterminado{}\% del corpus corresponde a mensajes desprovistos de carga afectiva explícita. Este resultado no constituye una deficiencia del motor léxico, sino que evidencia la función referencial-operativa predominante en la comunicación universitaria. La asignación formal de la etiqueta \emph{indeterminado} previene la inflación artificial de la categoría neutra.
  \item \textbf{Asimetría Afectiva Favorable:} Entre los mensajes con evidencia emocional suficiente, las expresiones positivas (\pctPositivos{}\%) duplican ampliamente a las negativas (\pctNegativos{}\%), manifestándose la \emph{\emocionTop{}} como el estado dominante, lo que confirma un ambiente de soporte y colaboración mutua entre estudiantes.
  \item \textbf{Estructura Temática y Dinámica Temporal:} El modelo LDA ($K=5$) agrupó de manera coherente las fases del semestre (trámites de Kardex, inscripciones presenciales, consultas de docentes y grupos de temporada). La actividad se sincroniza con el calendario institucional, alcanzando picos en \diaPico{} y en el rango de \horaPico{}.
  \item \textbf{Estabilidad Estocástica de Markov:} El análisis estocástico demostró que el estado neutro/factual actúa como un atractor del sistema ($\pi = 83,1\%$, permanencia media de \permanenciaNeutro{} turnos). Los episodios de frustración o estrés académico son transitorios y se disipan con resiliencia en un promedio de \recuperacionNegativa{} turnos.
  \item \textbf{Persistencia y Trazabilidad Transaccional CDC:} La migración hacia Neon Serverless PostgreSQL (16 tablas) y la implementación de triggers PL/pgSQL permitieron capturar cambios a nivel de atributo en JSONB, optimizándose la arquitectura mediante la amortización de triggers en tablas masivas temporales.
  \item \textbf{Paridad Matemática R vs. Python:} Se certificó la reproducibilidad científica estricta entre el entorno de investigación en R y el microservicio web en Flask, logrando un 100\% de éxito en la suite automatizada de pruebas unitarias (\texttt{pytest}) con tolerancia numérica nula ($\epsilon < 10^{-6}$).
\end{enumerate}

\subsection{Trabajo Futuro}

Como líneas de investigación y desarrollo futuro se plantean:
\begin{enumerate}
  \item \textbf{Modelos Híbridos Transformer en Español:} Incorporar modelos de lenguaje preentrenados ligeros (como BETO o RoBERTa-BNE) para procesar los mensajes catalogados como \emph{indeterminados}, reduciendo la brecha léxica hacia coberturas superiores al 85\% sin sacrificar la interpretabilidad.
  \item \textbf{Pool de Conexiones Multiplexado:} Implementar un pool de conexiones persistente en la capa Flask (\texttt{ThreadedConnectionPool}) para mitigar la latencia de handshakes SSL en bases de datos serverless.
  \item \textbf{Contenerización y Despliegue CI/CD:} Empaquetar la solución en contenedores Docker multi-stage para orquestar la ejecución de los pipelines y el dashboard en infraestructuras universitarias locales o en la nube.
\end{enumerate}
